\documentclass[./main.tex]{subfiles}
\begin{document}

% Slide # 4
\begin{frame}[t, label=slide04]
        % Title
        \frametitle{Taxonomy of tipping points}

        % Body
        \footnotesize

        Consider the following It\^{o} process

        \begin{equation*}
           \begin{cases}
               dx_{t} = f(x_{t},\mu(t))dt + \eta(x_{t})dW_{t}\,, \\
               \mu(t) = g(\varepsilon t)\,, 
           \end{cases}
        \end{equation*}
        
        where $f(x_{t},\mu)$ is the\alr{drift}(non-autonomous deterministic dynamics), $\eta(x_{t})$ is the\alr{diffusion}(stochastic fluctuations), $g(\varepsilon t)$ is the\alr{shift}(time-evolution of the parameter) and $W_{t}$ is Brownian motion.

        \vspace{3mm}

        Let's assume that $f(x;\mu)$ as at least one attractor $x_{\mu}$, in the frozen system at paramater value, $\mu$ with $\mathbb{B}(x_{\mu})$ being its basin of attraction.

        \uncover<2->{
                \begin{definition}[non-rigorous]\label{def:tipping_points}
                        We say that a sample path $x_{t} = x_{0} + \int_{0}^{t}\eta(x_{s})dW_{s} + \int_{0}^{t} f(x_{s},g(s))ds$ undergoes
                        \begin{itemize}
                                \uncover<3->{
                                        \item \textbf{R-tipping} if, in the deterministic limit $\eta\to0$, there is a finite $T>0$ s.t. $x_{T}\not\in \mathbb{B}(x_{\mu})$;
                                }
                                \uncover<4->{
                                        \item \textbf{N-tipping} if, in the singular limit $\varepsilon\to0$, there is a finite $T>0$ s.t. $x_{T}\not\in \mathbb{B}(x_{\mu})$;
                                }
                                \uncover<5->{
                                        \item \textbf{B-tipping} if there is a finite $T>0$ s.t. $x_{\mu}$ changes stability or no longer exists.
                                }
                        \end{itemize}
                \end{definition}
        }

\end{frame}

\end{document}
